\documentclass{article}

\title{Ch. 25 CW}
\author{Principles of Macroeconomics}
\date{\today}

\begin{document}

\maketitle

\begin{enumerate}

\item Government survey takers determie that typical family expenditures each month in the year designated as the base year are as follows;
	\begin{center}
	\begin{tabular}{c c}
	25 pizzas & \$10 each \\
	Apartment rent & \$600 per month \\
	Gasoline \& car maintenance & \$100 \\
	Cell phone service & \$50
	\end{tabular}
	\end{center}
In the year following the base year, the survey takers determine that pizzas have risen to \$11 each, apartment rent has risen to \$700 per month, gasoline and maintenance have risen to \$120, and phone service has dropped to \$40.

	\begin{enumerate}

	\item Find the CPI in the subsequent year. 
	
	\item Find the rate of inflation between the base year and the subsequent year.
	
	\item If a family's nominal income rose by 5\% between the base year and the subsequent year, are they worse off or better off in terms of what their income is able to buy?
	
	\end{enumerate}

\item According to the U.S. Census Bureau, nominal income for the median family of four in the United States was \$23,618 in 1985, \$34,076 in 1995, \$46,326 in 2005, and \$57,230 in 2015. In purchasing power terms, how did family income compare in each of those four years? The CPI (1983 = 100) was 107.6 in 1985, 152.4 in 1995, 195.3 in 2005, and 237 in 2015. In which decades did the median family's purchasing power increase?

\item The typical consumer's food basket in the base year 2022 is as follows:
	\begin{center}
	\begin{tabular}{c}
	30 chickens at \$4 each \\
	10 hams at \$8 each \\
	10 steaks at \$8 each
	\end{tabular}
	\end{center}
In 2023, the price of chicken remains the same; the price of ham rises to \$9, and the price of steak rises to \$10.

	\begin{enumerate}
	
	\item Calculate the percent change in the "cost-of-eating" index between 2022 and 2023.
	
	\item Suppose that consumers are completely indifferent between two chickens and one ham. How large is the substitution bias in your answer to part (a)?
	
	\end{enumerate}
	
\item On January 1, 2022, Albert invested \$1,000 at 6\% interest per year for three years. On January 1, 2022, the CPI was 100. On January 1, 2025, it was 118. Albert reinvests interest earnings each year. What was the percent change in Albert's purchasing power from 2022 to 2025?

\item Frank is lending \$1,000 to Sarah for two years. Frank and Sarah agree that Frank should earn a 2\% real interest rate per year. The CPI is 100 when Frank makes the loan. They expect it to be 110 in one year and 121 in two years. What nominal interest rate should Frank charge Sarah?

\item Rosalie the Retiree knows that when she retires in 16 years, her company will give her a one-time payment of \$20,000. However, if the inflation rate is 6\% per year, how much buying power will that \$20,000 have when measured in today's dollars?

\item In 2003, the commentator Jim Cramer declared, "Your home, alas, is not a dot-com, soon to be deflated and delisted. It's simply the best investment you'll ever make." While this advice seems dubious in light of the housing collapse of 2008, perhaps it makes sense when we take a longer-run perspective. The median sale price of a home in 1963 was \$18,000, while the median sale price of a home in 2025 was \$426,800. In 1963, the CPI was 30.6, and in 2025, it was 322. 

Compute the average annual real rate of return on housing from 1963 to 2025. The real rate of return is the percent growth in value. How does this compare to 3.5\%, the average real rate of return on long-term Treasury bonds (notoriously safe, but low return assets)? 

\end{enumerate}

\end{document}